\section{Prolegomena to Historiography}

As an experiment in historiography, I was planning to analyze the recent election in the USA. However, it became obvious that some prefatory notes are necessary. It is pointless to repeat what has already been said ten thousand times, words that convince no opponent and shed little light for their supporters because they are based on unsupported opinion rather than metaphysical principles. We take it for granted, at least in small things, that acting against reality will lead to dissatisfactory results. At the physical level, the negative outcome may be immediate: a child learns that falling or touching a hot object leads to immediate pain.

Unfortunately, the feedback may not be immediate. For example, a woman may go from one abusive relationship to the next before she understands, if ever, the role she herself was playing. At the social level, it may take generations for the lesson to be learned. As it is said, the iniquities of the fathers fall on the children. Very few men have the acuity to see how decisions made today will affect their descendants three generations hence. Fundamentally, this is because the deciders themselves are ignorant of the cosmic Order and act against it. They seldom experience the consequences of their acts.

So in the midst of this endless discussion, no one ever bothers to pause to ask why people hold the opinions they do, why they persist in holding them in the face of all illogic and counterfactual evidence, and why all debate ends in aporia rather than a higher resolution. If we start with the reasonable assumption that all rational men act in accordance to what they consider to be real, the only conclusion is that men have different notions of what is real. The remarkable fact is that men will not budge from their positions. No matter how poorly he may have done in school or on tests—which should certainly bring a measure of intellectual humility—every man is convinced he understands the real.

For a few, there is an understanding of the transcendent as real. For others, it is nobility and virtue. A much larger group can only accept the material world as real, while the masses have no reality beyond their own satisfactions. These various understandings of the real are characteristic of caste. It is important to keep in mind that caste is a metaphysical principle, i.e., an expression of the very nature of things, hence it surpasses any moral or psychological considerations. Specifically, caste is not detectable by psychological methods, nor is it ``curable". Furthermore, there is no moral significance to caste, so any discussion of their characteristics does not imply any moral judgment, whether positive or negative.

Nevertheless, the same cannot be said for any disorder in their power relationships. In a Traditional society those who know, will lead, and the rest will believe depending on their positions. For example, the man whose concern is with the physical world will nevertheless be able to participate in the transcendent through his faith in his leaders. Of course, the modern world rejects this, because of its commitment to egalitarianism; all men are equal insofar as they are animals, but not in respect to their Intellects. Only in the primordial state are there no caste distinctions.

Fundamental to this historiography is the notion of Order in general, and specifically how that manifests in the relationship of the different castes to each other. Although we have written enough posts about the castes of those who pray (Brahman) and those who fight (Kshatriya), their influence on the modern world is essentially negligible. Since he rule of the modern world is in the hands of those who work, the characteristics of the relevant castes need to be more fully fleshed out.

Frithjof Schuon's discussion of castes provides a fuller discussion of this topic. The castes categorize the different way that men understand what is real.

\begin{itemize}
\item \textbf{Brahman}. Those who pray. This is the intellectual, contemplative and sacerdotal type. For him, the changeless and the transcendent are real. He is an idealist (i.e., ideas are real) and objective (he can transcend his own perspective). 
\item \textbf{Kshatriya}. Those who fight. This is the warrior or the administrator type. For him, it is action that is real. His aggression is tempered by generosity, and his passionate nature by nobility, self-control, and greatness of soul. He values action, honour, virtue, glory and nobility. This caste is also idealist, but in a subjective way. 
\item \textbf{Vaishya (or bourgeois)}. This is the first division of those who work. He is the merchant, the peasant, the artisan, the man whose activities are directly bound up with material values. For him, it is riches, security, prosperity and well-being that are real. He even thinks of religion in a material and economic way, which he understands as the need to accumulate merit. This caste is materialist and objective. 
\end{itemize}
These three castes constitute the twice-born, or spirits with a body. Beyond them, men can be described only as a body with a human consciousness.

\begin{itemize}
\item \textbf{Shudra (or serf)}. Of those who work, the shudra who can only perform manual work of a quantitative type. For him, it is bodily things that are real, such as eating, drinking, the satisfaction of physical needs, and all other corporeal or sensual activities. This caste is materialist and subjective. Thus, the shudra is passive in relation to matter, hence he cannot govern himself. Therefore, he is dependent on the will of another which, in a well-ordered society, would be the spiritual authorities who guide and protect him. Their religion is the one that makes them feel good. 
\item \textbf{Chandala}. This is the outcast, the man lacking caste because of his chaotic character. He exists because he realizes the possibilities that other men reject and exists on the fringe of a normal society. Schuon claims he may appear as a ``chimney sweep, acrobat, comedian, or executioner, not to mention illicit occupations." 
\end{itemize}
Those who have understood thus far will come to some immediate conclusions. The USA is under the dominance of the bourgeois, serfs, and even outcasts. It is their rivalries and power struggles that determine events. There is no possibility of returning to ``traditional values" in this milieu. Just as the woodsmen can determine what animals surround him at night in the woods just by listening carefully to their sounds, so can we hear the words rise up from the various castes. Those of the bourgeoisie insist that ``it is the economy, stupid." The serfs are interested in physical satisfactions, hence they want to keep feticide, sodomy, and drug use or drunkenness (meant in the larger sense as the deliberately induced impairment of normal conscious activities) all legal. Even the outcasts have a place in this society. An entire generation of American youths get their news from a comedian. Rock stars and actors mould public opinion. The USA admires in a crude way the gangster, the prostitute, the pusher, and so on. The ugly, the dissonant, the unusual, the freakish, the perverted have destroyed their taste for the beautiful, the harmonious, the normal, the sane, the healthy.

There are some direct quotations from Castes and Races by Frithjof Schuon.



\flrightit{Posted on 2012-11-13 by Cologero }

\begin{center}* * *\end{center}

\begin{footnotesize}\begin{sffamily}



\texttt{Avery on 2012-11-13 at 22:14 said: }

``A much larger group can only accept the material world as real, while the masses have no reality beyond their own satisfactions."

This is an important distinction and worth explaining in full. After all, we live in an age where such an esteemed moral philosopher as Richard Dawkins conceives of reality as a ``virtual world" created inside our heads, without any understanding of the values that are caused by such a view.


\hfill

\texttt{Cologero on 2012-11-13 at 23:47 said: }

Can we conclude that Dawkins, then, is a subjective materialist, i.e., a serf? Since his mentality is limited, you can see there is no fact or logical argument to budge him from his conception of reality as mind produced. Given his assumptions, has he ever been able to explain how people of quite modest intelligence are able to create such complex ``virtual worlds"?


\hfill

\texttt{Aziz Ali on 2012-11-14 at 05:27 said: }

How does one find out which caste one belongs to? On one hand I have an interest in this kind of ``unconventional" spirituality/metaphysics, on the other hand I do enjoy standup-comedy and hip-hop \& r `n b (albeit not exclusively). If I am a shudra or a tschandala, maybe I'm missing the point of reading these kinds of blogs and related literature and reading them for the wrong reasons. Isn't the most Traditional thing of me to do, then, to stop caring about Tradition and do what I'm metaphysically designed to do best?


\hfill

\texttt{Cologero on 2012-11-14 at 07:49 said: }

If only life were so easy, that we could categorize everyone, put him in a box, and arrange him on a shelf. The Spirit blows where He wills and we judge not by appearances. There have been crazy saints and men of blame (malamatiyah).

Schuon aptly puts it:

\begin{quotex}
If religious hypocrisy is an inevitable fact, the contrary must also be possible, namely wisdom and virtue hiding under appearances of scandal.

\end{quotex}
As Guenon points out, today hardly anyone is in his natural place. We judge caste by level of understanding; a shudra personality would have little interest in anything truly metaphysical. At the level of knowledge, we are the same, but that does not imply outer conformity.

Read what Nietzsche writes about giving style to one's character\footnote{\url{http://www.gornahoor.net/?p=11}} and then avoid being gloomy.


\hfill

\texttt{Avery Morrow on 2012-11-14 at 08:27 said: }

The list above is intended to answer that question, but if you don't recognize yourself in those spiritual categories, let me restate them a little:

Brahmans have direct access to the Divine. They understand their life calling and can intuit traditional knowledge. A brahman wouldn't be offended by stand-up comedy, understanding the nature of its arising, but wouldn't need it to enjoy life either.

Kshatriyas have knowledge of virtue and, if they understand it properly, can attribute virtue to the Divine. Kshatriyas become fed up with a world full of heterodoxies and settle into a single orthodoxy. They are the caste you may be imagining who would be offended by rap music if they are not members of the culture that actually produces rap.

Vaishyas, when directed properly, believe in reality and therefore in an Order to reality. When misdirected by outside forces, they stop believing in order and pretend to deny reality (Dawkins falls into this category), and instead propose that traditional institutions should revise their principles to meet the needs of the lower castes, and that all metaphysics is simply an expression of taste, preference, and sentiment.

Shudras simply don't believe that anything can exist outside their perceptions. They believe that religion is only good for getting what you want. They have no conception of metaphysics and would not understand any of this blog. Dawkins does not actually fall into this category because if he did, he wouldn't be writing books, he would just be hanging out in pubs using his high-falutin fancy talk to scam men and seduce women.

Outcasts, for some reason or another, have been excluded from mainstream society so they lack access to their society's tradition, including what Guenon calls the ``initiation". Instead, the force that dominates their minds is chaos, and they must take whatever work they can get in a world of chaos. An outcast can maneuver his way into one society or another if he works hard for it, but he can also be manipulated as a disruptive agent if he does not try to deepen his understanding.

Note that Western society, i.e. late modernity, attempts to make life easy for shudras and outcasts especially.


\hfill

\texttt{Avery Morrow on 2012-11-14 at 08:37 said: }

By the way, that comment was completely off the top of my head, as part of the conversation about this post; I haven't read as much as Cologero on this subject and I don't mean to contradict any information he has.


\hfill

\texttt{Logres on 2012-11-14 at 09:07 said: }

Just from watching the bookshelves, you can see that most (or at least a large part) of American new literature is about shudras and outcastes – that is, goldenhearted crazy people who ``make good" in the end, one way or the other. TV, on the other hand, seems to still revolve around the merchant caste, mainly. Is this anyone else's experience?


\hfill

\texttt{Jason-Adam on 2012-11-14 at 17:14 said: }

From my own readings, I've come to feel that I am a debased Kshatriya, debased because I am not able to perform my natural function in this society. One aspect of my Kshatriya nature is that I'm not interested so much in Tradition as an abstract ideal but rather Tradition as the knowledge and practice of my ancestors. I love to learn but at the same time find learning useless unless it is combined with actual actions that I can perform.To the modern world, I am an outcast because I find that my personality is too strong to fit in with the mediocrity and nihilism of the masses. 

I think there is a problem that Kshatriyas, if they remain uninitiated, can find themselves debased by the modern world into chandalas and warriors for the forces of destruction.


\hfill

\texttt{Michael on 2012-11-14 at 18:13 said: }

Is there a connection between caste and European nobility? Is it fair to assume that a European noble belongs to one of the upper castes?


\hfill

\texttt{Jason-Adam on 2012-11-14 at 18:23 said: }

Used to be the case but now I question it – I've known some people of noble descent who behaved no better than ghetto scum….


\hfill

\texttt{Dominion on 2012-11-14 at 19:24 said: }

I'd say this was probably the case in the first few generations, when the caste was proven through warfare and when new blood could enter it at will so long as they showed the characteristics of the caste. When it became institutionalized and `closed', the degeneration began.


\hfill

\texttt{Saladin on 2012-11-14 at 21:45 said: }

Excellent observations! I agree with almost all of the ideas expressed except that ``the rule of the modern world is in the hands of those who work [i.e. Vaishias and Shudras]." I disagree! I am certain that you are well aware of `La Guerre occulte' waged against every Tradition on the face of the earth. This occult war is not directed by ``those who work". These multitudes only serve as tools or at best auxiliaries of the real masters of the Anti-traditional party. Therefore, if we only consider the human ``leaders" of the Counter-Initiation (and leave the more sinister forces and entities who control these men), we are forced to acknowledge that the masters of the modern world are either castless or Brahmans who have gone astray (if such a thing is in fact possible).


\hfill

\texttt{Mihai on 2012-11-15 at 04:07 said: }

Dawkins never used the term ``moral philosopher" to describe himself, nor any kind of ``philosopher" for that matter. He writes about religious and/or philosophical matters because he believes that anyone can talk about these things according to fancy, since they do not come under any empirical observation. 

His works on religion and morals display a complete ignorance of even high-school level philosophy and logic. In general they are mere rambling, without any basic knowledge about the topic discussed in them.


\hfill

\texttt{Avery Morrow on 2012-11-16 at 22:44 said: }

Yes, I was being sarcastic.


\hfill

\texttt{Max on 2015-05-02 at 10:29 said: }

The third caste have been on the decline and is becoming rarer. How often today is someone actually producing something of quality, just speaking in material terms? Capitalism is likewise uncommon even though some cling on to names without understanding their meaning. Enlightenment-type rationality is almost gone, with universities now occupied with conforming to popular opinion. Shudras are a majority, but what is more disturbing is an idealization of outcasts. They have been moving behind the scenes, manipulatively providing entertainment or whatever for the masses. The bourgeois would have found that repulsive, but they have been replaced by shudras on many official positions.


\end{sffamily}\end{footnotesize}
